% Generated draft: requires mathematical and source review.
\documentclass[11pt]{article}
\usepackage[margin=25mm]{geometry}
\usepackage{fontspec}
\setmainfont{Latin Modern Roman}
\usepackage{amsmath,amssymb,mathtools,unicode-math}
\setmathfont{Latin Modern Math}
\usepackage{enumitem,needspace,xurl,hyperref}
\setlength{\parindent}{0pt}
\setlength{\parskip}{5pt}
\setlength{\emergencystretch}{3em}
\newcommand{\N}{\mathbb{N}}
\newcommand{\Z}{\mathbb{Z}}
\newcommand{\Q}{\mathbb{Q}}
\newcommand{\R}{\mathbb{R}}
\newcommand{\C}{\mathbb{C}}
\newcommand{\F}{\mathbb{F}}
\newcommand{\A}{\mathbb{A}}
\newcommand{\PP}{\mathbb{P}}
\newcommand{\E}{\mathbb{E}}
\newcommand{\eps}{\varepsilon}
\newcommand{\dd}{\,\mathrm{d}}
\begin{document}
\section*{UnsolvedMath ID: 30004980}
\textbf{Source problem code:} catalogue entry
\subsection{Definitions and mathematical statement}
Under the conventional interpretation of ``graph'' as a finite simple undirected graph, define twin-width as follows. This graph convention is not explicitly specified in the supplied source record.

Let \(G=(V,E)\) have \(n\geq 1\) vertices. A contraction sequence is a sequence \(\mathcal P_n,\mathcal P_{n-1},\ldots,\mathcal P_1\) of partitions of \(V\), where \(\mathcal P_n=\{\{v\}:v\in V\}\), \(\mathcal P_1=\{V\}\), and each \(\mathcal P_{i-1}\) is obtained from \(\mathcal P_i\) by merging two parts. For distinct parts \(A,B\in\mathcal P_i\), call their relation red if there exist \(a,a'\in A\) and \(b,b'\in B\) such that \(\{a,b\}\in E\) and \(\{a',b'\}\notin E\). Thus the adjacency relation between \(A\) and \(B\) is neither complete nor empty. Set
\[
 r_{\mathcal P_i}(A)=\bigl|\{B\in\mathcal P_i\setminus\{A\}:\text{the relation between }A\text{ and }B\text{ is red}\}\bigr|.
\]
The twin-width of \(G\) is
\[
 \operatorname{tww}(G)=\min_{(\mathcal P_n,\ldots,\mathcal P_1)}\ \max_{1\leq i\leq n}\ \max_{A\in\mathcal P_i} r_{\mathcal P_i}(A).
\]
For every positive integer \(n\), define
\[
 M(n)=\max\{\operatorname{tww}(G):G\text{ is a finite simple undirected graph with }|V(G)|=n\}.
\]
The problem is to determine \(M(n)\) for every positive integer \(n\).

The supplied dated assessment describes \((n-1)/2\) as the conjectured exact maximum. Its intended quantifiers, admissible orders, or rounding convention are unresolved: \(M(n)\) is integer-valued, whereas \((n-1)/2\) is not an integer for even \(n\). Consequently, that assessment is retained as supplied context, rather than imposed as an unconditional equality for every \(n\). The source does not specify whether the empty graph order \(n=0\) is intended.

The raw original-statement field continues into a separate problem on treewidth of hereditary graph classes and ends in a truncated variant. That continuation is recorded in the accompanying notes and references; its hypotheses do not occur in the supplied clean twin-width question.
\subsection{Short English statement}
For each positive integer \(n\), determine the largest twin-width of an \(n\)-vertex graph, using the conventional finite simple undirected graph interpretation. Twin-width measures the smallest possible maximum number of mixed adjacency relations encountered while repeatedly merging vertex parts. The supplied conjectural expression \((n-1)/2\) needs clarification about its intended orders or rounding.
\subsection{Sources}
\begin{itemize}
\item Local UnsolvedMath registry, permanent ID \(30004980\), title ``Maximum Twin-Width of \(n\)-Vertex Graphs,'' source problem code \texttt{catalogue entry}, in the Oberwolfach Reports Open Problems collection. The record lists proposed year 2022 and supplies no proposer attribution for the twin-width question.
\item Supplied source citation: \emph{Graph Theory (2023)}, Oberwolfach Reports, DOI 10.4171/owr/2022/1, pp. 5--78. \url{https://doi.org/10.4171/owr/2022/1}. The record also labels this reference ``Oberwolfach Report 1/2022.'' The discrepancy between the supplied title year and report designation is unresolved.
\item Supplied literature source: Ahn, Hendrey, Kim, and Oum, ``bounds for twin-width.'' \url{https://arxiv.org/abs/2110.03957}.
\item Supplied literature source: ``Oberwolfach Report 1/2022.'' \url{https://doi.org/10.4171/owr/2022/1}.
\item Literature assessment dated exactly \texttt{2026-08-22}: the supplied curation labels the problem ``partially solved'' and states that general \(n\)-vertex graphs have twin-width at most \(n/2\) plus a lower-order term, that conference graphs attain \((n-1)/2\), and that the conjectured exact maximum \((n-1)/2\) remains open. This is supplied dated curation, not an independently verified literature assessment. The record itself requests independent verification of time-sensitive status and cited claims.
\item References embedded in the separate treewidth continuation: Tara Abrishami, Maria Chudnovsky, Sepehr Hajebi, and Sophie Spirkl, \emph{Induced subgraphs and tree-decompositions III. Three-path-configurations and logarithmic tree-width}, supplied identifier arXiv:2109.01310; and Ni Luh Dewi Sintiari and Nicolas Trotignon, \emph{(Theta, triangle)-free and (even hole, K4)-free graphs. Part 1: Layered wheels}, supplied identifier arXiv:1906.10998. These references belong to the appended treewidth discussion, not to an established attribution of the twin-width problem.
\end{itemize}
\end{document}
